% ============================================================
% LAMPIRAN A - KATALOG PERINTAH VERIFIKASI USE CASE
% Katalog perintah yang dapat dijalankan ulang (read-only terhadap
% kluster langsung) untuk membuktikan tiap skenario SK-F/SK-N pada
% Bab VI. Dikelompokkan mengikuti tujuan platform T-1..T-4 agar rantai
% T-N <-> SK <-> KF/KNF tetap satu lawan satu. Bersumber dari direktori
% experiments/ (sejajar platform/ dan use-case-crypto/); rincian argumen
% tiap skenario ada pada README.md di setiap subdirektori.
% ============================================================
\cleardoublepage
\chapter{KATALOG PERINTAH VERIFIKASI USE CASE}
\label{appx:verifikasi}

Lampiran ini memuat katalog perintah yang dapat dijalankan ulang untuk membuktikan jalur ujung ke ujung pada \textit{use case} pasar aset kripto. Setiap perintah dijalankan secara manual dan bersifat \textit{read-only} terhadap kluster yang sedang berjalan, sehingga tidak menjadi bagian dari \texttt{make phase-full} maupun rekonsiliasi Argo CD dan tidak mengubah sistem yang diuji, sejalan dengan metode evaluasi pada Bab~\ref{chap:evaluasi}. Berkas skenario berada pada direktori \texttt{experiments/} yang sejajar dengan \texttt{platform/} dan \texttt{use-case-crypto/}; rincian argumen tiap skenario tersedia pada berkas \texttt{README.md} di masing-masing subdirektori. Prasyarat eksekusi adalah platform telah terpasang (\texttt{make phase-full}), \textit{use case} telah ter-\textit{deploy} (\texttt{make usecase-crypto-build \&\& make usecase-crypto-up}), dan seluruh \textit{pod} berstatus \textit{Running}. Sebelum menjalankan skenario apa pun, konfigurasi domain dimuat dan \textit{namespace} pembangkit beban dibuat melalui preambel berikut.

\begin{lstlisting}[caption={Preambel: pemuatan konfigurasi domain dan \textit{namespace} pembangkit beban},basicstyle=\ttfamily\footnotesize,breaklines=true,captionpos=t]
# Muat blok DOMAIN dari satu konfigurasi (entitas, tabel medali, endpoint,
# knob beban). Ganti use case cukup ganti berkas ini, bukan kode skenario.
set -a; . "${EXPERIMENTS_CONFIG:-experiments/config.crypto.env}"; set +a

# Namespace pembangkit beban (k6, chaos), terpisah dari produksi.
kubectl apply -f experiments/namespace.yaml
\end{lstlisting}

\section{Tujuan T-1: Arsitektur Terintegrasi dan GitOps}
\label{sec:appx-t1}

Kelompok T-1 membuktikan pemasangan deklaratif dari satu sumber Git, otomasi siklus hidup model setara MLOps Level 2, serta atribut kualitas lintas komponen. Perintahnya dirangkum pada Listing~\ref{lst:cat-t1}.

\begin{lstlisting}[caption={Perintah verifikasi tujuan T-1},label={lst:cat-t1},basicstyle=\ttfamily\footnotesize,breaklines=true,captionpos=t]
# SK-F-01 (KF-01): rekonsiliasi GitOps feature view via Argo CD. Definisi feature
# view Feast ada pada ConfigMap feast-feature-repo (overlay local menambah prefix
# crypto-, menjadi crypto-feast-feature-repo); use-case-crypto repo terlacak Gitea.
$EDITOR use-case-crypto/manifests/base/configmaps/feast.yaml
git -C use-case-crypto commit -am "feat: feature view" && git -C use-case-crypto push gitea
argocd app sync crypto-use-case-local && argocd app diff crypto-use-case-local
# Definisi tersusun menjadi repo Feast lalu diregistrasi task materialize (Airflow)
# via 'feast apply' tanpa langkah manual; 'feast feature-views list' dari pod
# ber-repo Feast menampilkannya. Konfirmasi reconcile read-only di kluster:
kubectl -n use-case-crypto get configmap crypto-feast-feature-repo \
  -o jsonpath='{.data.definitions\.py}' | grep 'FeatureView('

# SK-F-12 (KF-12): konfigurasi deklaratif, konvergen tanpa drift.
argocd app diff crypto-use-case-local
kubectl -n use-case-crypto get deploy

# SK-F-06 (KF-06): promosi kanari dan rollback via Argo Rollouts.
kubectl -n use-case-crypto get rollouts
kubectl argo rollouts -n use-case-crypto get rollout <rollout> --watch   # promote | abort

# SK-F-13 (KF-13): notebook satu environment uv memuat SDK Feast, MLflow, dan
# klien KServe v2, lalu menjalankan ketiganya sebagai operasi baca independen.
# Dijalankan dari Kubeflow Notebook in-cluster tempat repo Feast dan DNS layanan
# tersedia. Verifikasi struktural: environment terpasang, notebook tereksekusi.
cd use-case-crypto/notebooks
uv sync
uv run jupyter nbconvert --to notebook --execute unified_sdk_walkthrough.ipynb
kubectl get inferenceservice -A | grep predictor   # endpoint v2 aktif (crypto-predictor)

# SK-F-14 (KF-14): jenuhkan ClusterQueue Kueue, job luar kuota Pending.
kubectl get clusterqueue,localqueue,workloads -A

# SK-N-03 (KNF-03): skalabilitas horizontal HPA dan KEDA.
k6 run experiments/load/hpa-keda-rampup.js
kubectl -n use-case-crypto get hpa,scaledobject,pods -w

# SK-N-04 (KNF-04): injeksi fault pod dan jaringan (RTO<5m, RPO<1m).
envsubst < experiments/chaos/pod-chaos.yaml | kubectl apply -f -
kubectl apply -f experiments/chaos/network-chaos.yaml

# SK-N-08 (KNF-08): ganti runtime serving via CRD dan overlay Kustomize.
kubectl apply -k use-case-crypto/manifests/overlays/local

# SK-N-09 (KNF-09): SSO satu sesi dex dan oauth2-proxy.
kubectl get pods -A | grep -E 'dex|oauth2-proxy'

# SK-N-10 (KNF-10): rasio kompresi dan atribusi biaya OpenCost.
kubectl get pods -A | grep -i opencost

# SK-N-11 (KNF-11): rekonstruksi penuh node k3s baru dari Git.
make phase-full

# SK-N-12 (KNF-12): cakupan uji dan waktu konvergensi phase-full.
make usecase-crypto-test
time make phase-full
\end{lstlisting}

\section{Tujuan T-2: Tata Kelola Data}
\label{sec:appx-t2}

Kelompok T-2 membuktikan katalog metadata, \textit{lineage}, kontrak kualitas, observabilitas, dan keamanan berjalan otomatis sepanjang \textit{pipeline}. Perintahnya dirangkum pada Listing~\ref{lst:cat-t2}.

\begin{lstlisting}[caption={Perintah verifikasi tujuan T-2},label={lst:cat-t2},basicstyle=\ttfamily\footnotesize,breaklines=true,captionpos=t]
# SK-F-08 (KF-08): lineage end-to-end DataHub, sumber hingga prediksi.
experiments/lineage/lineage-walk.sh discover
experiments/lineage/lineage-walk.sh walk '<urn>'         # urn dari hasil discover

# SK-F-09 (KF-09): dua run pelatihan, replay snapshot, delta metrik di bawah 1%.
kubectl get pods -A | grep -i mlflow                     # konsol MLflow via ingress

# SK-F-10 (KF-10): transisi status model dan audit trail MLflow Registry.
kubectl get pods -A | grep -i mlflow

# SK-N-05 (KNF-05): checkpoint Great Expectations menolak pelanggaran kontrak.
kubectl -n use-case-crypto get cronjob | grep -iE 'quality|expect|ge-'
kubectl -n use-case-crypto create job ge-check --from=cronjob/<ge-checkpoint>

# SK-N-06 (KNF-06): telusuri satu trace OpenTelemetry di Tempo via Grafana.
kubectl get pods -A | grep -E 'tempo|grafana'

# SK-N-07 (KNF-07): pindai Trivy dan periksa TLS 1.3 serta AES-256.
kubectl get vulnerabilityreports -A
testssl.sh <public-endpoint>
\end{lstlisting}

\section{Tujuan T-3: Deteksi \textit{Drift} dan Pelatihan Ulang}
\label{sec:appx-t3}

Kelompok T-3 membuktikan deteksi \textit{drift} terotomasi memicu pelatihan ulang dan penyusunan fitur pelatihan menjaga \textit{point-in-time correctness}. Perintahnya dirangkum pada Listing~\ref{lst:cat-t3}.

\begin{lstlisting}[caption={Perintah verifikasi tujuan T-3},label={lst:cat-t3},basicstyle=\ttfamily\footnotesize,breaklines=true,captionpos=t]
# SK-F-03 (KF-03): point-in-time correctness. Probe PIT manual via SDK Feast
# (prosedur Bab VI metode T-3); perintah ini hanya memastikan substrat gold mendarat.
# Kredensial admin ClickHouse diambil dari Secret bila user in-pod bersandi.
export CLICKHOUSE_USER=$(kubectl -n "$CLICKHOUSE_NAMESPACE" get secret clickhouse-admin \
  -o jsonpath='{.data.CLICKHOUSE_USER}' | base64 -d)
export CLICKHOUSE_PASSWORD=$(kubectl -n "$CLICKHOUSE_NAMESPACE" get secret clickhouse-admin \
  -o jsonpath='{.data.CLICKHOUSE_PASSWORD}' | base64 -d)
experiments/etl/data-landed-check.sh

# SK-F-05 (KF-05): otomasi pelatihan Kubeflow, baca Feast tulis MLflow.
uv run --with 'kfp[kubernetes]==2.16.0' use-case-crypto/pipelines/retraining_pipeline.py

# SK-F-07 (KF-07): geser distribusi, picu detektor drift dan retrain-on-drift.
uv run experiments/drift/inject_drift.py --sigma 2.0
kubectl -n use-case-crypto get cronworkflow,workflow
\end{lstlisting}

\section{Tujuan T-4: Layanan Fitur \textit{Dual-Store}}
\label{sec:appx-t4}

Kelompok T-4 membuktikan layanan fitur \textit{dual-store} menyajikan fitur tabular, agregat, dan vektor dengan kontrak konsisten antara pelatihan dan penyajian, beserta latensi dan \textit{throughput} jalur penyajian. Perintahnya dirangkum pada Listing~\ref{lst:cat-t4}.

\begin{lstlisting}[caption={Perintah verifikasi tujuan T-4},label={lst:cat-t4},basicstyle=\ttfamily\footnotesize,breaklines=true,captionpos=t]
# SK-F-02 (KF-02): online non-null Feast dan substrat offline ClickHouse.
kubectl -n model-lifecycle port-forward svc/feast 6566:6566 &
FEAST_URL=http://localhost:6566 uv run experiments/feast/online-serving-check.py
# substrat offline (ekspor kredensial ClickHouse seperti Listing T-3 bila bersandi):
experiments/etl/data-landed-check.sh

# SK-F-04 (KF-04): pencarian vektor Qdrant, recall dan latensi.
k6 run experiments/load/vector-search-latency.js

# SK-F-11 (KF-11): paritas batch (dbt volume_sma_20) vs stream (Flink secondary_avg).
# Cek substrat dulu; kredensial ClickHouse ekspor seperti Listing T-3 bila bersandi.
experiments/etl/data-landed-check.sh
experiments/parity/parity-check.sh

# SK-N-01 (KNF-01): latensi feature serving dan vector search (p99 sub-detik).
k6 run experiments/load/feast-online-latency.js

# SK-N-02 (KNF-02): throughput serving dan stream, pantau consumer lag.
k6 run experiments/load/feast-throughput.js
# Sisi stream: kredensial SASL Kafka dari kluster (tidak dikomit).
# Topik crypto.validated; pantau metrik kafka_consumergroup_lag.
kafka-producer-perf-test --topic crypto.validated --num-records 1000000 \
  --record-size 256 --throughput -1 \
  --producer-props bootstrap.servers="$KAFKA_BOOTSTRAP" \
  --producer.config <sasl.properties>
\end{lstlisting}

Setiap perintah pada katalog ini berpasangan satu lawan satu dengan baris Tabel~\ref{tab:uji_penerimaan} dan Tabel~\ref{tab:evaluasi_target} pada Bab~\ref{chap:evaluasi}, sehingga klaim pemenuhan kebutuhan dapat ditelusuri langsung ke perintah yang menghasilkan buktinya. Untuk \textit{use case} lain, hanya berkas \texttt{config.crypto.env} yang disalin dan disesuaikan blok DOMAIN-nya; kode skenario tetap \textit{domain-agnostic} dan tidak perlu disentuh.
